\section{Version 2: Path Coloring} \label{sec:kps_v2}

To extract a heavier set of paths from the cycle cover, the second iteration of
the generic algorithm utilizes non-cycle edges. To establish the core mechanics,
we first present a simplified variant before introducing the fully optimized
version that achieves the stronger approximation bound.

\subsection{The Simple Variant of Version 2}
The basic variant operates through the following sequence of
steps:\begin{enumerate}[label=\textbf{Step 2\alph*.}, leftmargin=*]
    \item Form an auxiliary undirected graph $G'$ derived from the base graph
    $G$ and its cycle cover $\mathcal{C}$ in the following manner:
    \begin{itemize}
        \item Substitute each 2-cycle $\{(u, v), (v, u)\}$ of $\mathcal{C}$ with
        an undirected edge $\{u, v\}$ carrying a weight of $0$.
        \item For all remaining pairs of vertices $u$ and $v$, introduce an
        undirected edge $\{u, v\}$ whose weight is equal to $\max\{w(u, v), w(v,
        u)\}$.
    \end{itemize}
    \item Compute a maximum-weight matching $M$ within this auxiliary graph
    $G'$.
    \item Revert the edges of $M$ to their directed forms, and augment this set
    by selecting a single edge from every $3^+$-cycle (ensuring no cycles are
    created in the corresponding undirected graph).\footnote{Proof that it can
    be done is omitted, because in the real variant we use the stronger Lemma
    \ref{lemma:path_coloring}.}
    \item Contract the 2-cycles into single vertices.
    \item Two-color edges of the remaining graph to divide it into two sets of
    directed paths. Form the final path collection $P$ by taking the edges from
    the heavier color group and supplementing them with one compatible edge
    extracted from each 2-cycle.
\end{enumerate}

\paragraph{Analysis of the Simple Variant}
Let $T_{opt}$ represent the optimal maximum-weight tour, with $w(T_{opt})$
denoting its total weight. When evaluating the weight of $T_{opt}$ within the
context of $G'$, we note that 2-cycles are assigned a weight of 0, while all
other edges maintain or increase their weights. As a result, the weight of
$T_{opt}$ in $G'$ is bounded below by $w(T_{opt}) - bW$, given that each 2-cycle
can contribute at most one edge to $T_{opt}$.

Because $T_{opt}$ is a single cycle on an even number of vertices, its edges can
be alternately split into two matchings averaging at least
$\frac{1}{2}(w(T_{opt}) - bW)$. Therefore, the maximum-weight matching $M$
obtained in Step 2b is at least this average weight.

Adding one edge from each $3^+$-cycle in Step 2c contributes an additional
weight of $pW$. The total weight of this constructed set of edges is therefore
bounded below by:
\begin{equation}
\alpha = \frac{w(T_{opt}) - bW}{2} + pW
\end{equation}

Contracting the 2-cycles in Step 2d does not alter this total weight, as they
were initially zero-weighted in $G'$. We are then left with a collection of
paths and cycles. These edges can be two-colored (partitioned into two color
classes) such that each color class forms a collection of vertex-disjoint
directed paths. Consequently, the heavier of these two classes must contain a
total weight of at least $\frac{\alpha}{2}$.

Upon uncontracting the 2-cycles, it is mathematically guaranteed that we can
always select one edge per 2-cycle to add to our chosen color class while
maintaining a valid collection of disjoint paths. Assume, for the contradiction,
that this was impossible. This would imply that a 2-cycle has two incident edges
of the same color, both directed either towards or away from the cycle. However,
in the contracted graph, these two edges would have been directed either toward
or away from the contracted vertex—meaning they could not have validly received
the same color in a path coloring, establishing a contradiction.

This final addition gives a set of directed paths with a minimum weight of:
\begin{equation}
\frac{\alpha}{2} + cW = \frac{w(T_{opt}) - bW}{4} + \frac{pW}{2} + cW
\end{equation}
Substituting the bounds $p \le \frac{1}{3}(1 - b - c)$ and $W \ge w(T_{opt})$
into this equation results in a bound of $\frac{5}{4}(c + p)w(T_{opt})$.

By combining this bound with the $(1 - c - p)w(T_{opt})$ bound from Version 1
(Theorem~\ref{lemma:version1}) -- taking whichever of the two constructions
yields the heavier set of paths -- it guarantees a weight of at least:
\begin{equation}
\max\left\{(1 - c - p),\ \frac{5}{4}(c + p)\right\} w(T_{opt}).
\end{equation}
Writing $x = c + p$, the first term is decreasing in $x$ and the second is
increasing in $x$, so their maximum is minimized exactly where the two coincide:
\begin{equation}
1 - x = \frac{5}{4}x \iff x = \frac{4}{9},
\end{equation}
at which point both bounds equal $1 - x = \frac{5}{9}$. Hence, in the worst case
over all valid $b, c, p$, this basic variant is guaranteed to achieve a
$\frac{5}{9}$-approximation for the Max TSP.

\subsection{Version 2: The Real Variant}

To secure a stronger approximation bound, the simplified algorithm requires two major improvements. First, rather than assigning zero weight to the 2-cycle edges (which only guarantees the inclusion of the lighter edge), we allow a 2-cycle edge to be placed in the matching if the heavier edge significantly outweighs the lighter edge. To achieve this, we assign a weight of $2(b_i - c_i)$ to the undirected edge representing a 2-cycle $C_i$, where $b_i$ and $c_i$ are the weights of the heavy and light edges, respectively. Second, a specialized path-coloring lemma is utilized, permitting the inclusion of every edge in a $3^+$-cycle except the lightest.\\
\vspace{0.5em}
\noindent\textbf{Algorithm: Version 2, Real Variant}
\begin{enumerate}[label=\textbf{Step 2\alph*.}, leftmargin=*]
    \item Build the undirected graph $G'$ from $G$ and $\mathcal{C}$ with the
    following adjustments:
    \begin{itemize}
        \item Reclassify any 2-cycle where $b_i > 2c_i$ as a $3^+$-cycle.
        \item Replace each 2-cycle $\{(u, v), (v, u)\}$ of $\mathcal{C}$ with an
        undirected edge $\{u, v\}$ carrying a weight of $2(b_i - c_i)$.
        \item For all other pairs of vertices $u$ and $v$, create an undirected
        edge $\{u, v\}$ with weight $\max\{w(u, v), w(v, u)\}$.
    \end{itemize}
    \item Extract a maximum-weight matching $M$ from $G'$.
    \item To the directed versions of $M$'s edges, attach every edge besides
    lightest from the $3^+$-cycles.
    \item Contract the 2-cycles found in $\mathcal{C}$.
    \item Two-color edges of the remaining graph using Lemmas
    \ref{lemma:path_coloring} and \ref{lemma:induced_cycles} to generate two
    sets of directed paths. The final set $P$ comprises the heavier color group
    combined with one edge from each 2-cycle.
\end{enumerate}
\vspace{0.5em}

We start by stating the primary coloring lemma, referring to the resulting
coloring as a \textit{path coloring}. Proof of this was omitted in the original
paper, but since it is a complex result, we devote Chapter 4 to its elaboration.

\begin{lemma}[Path-Coloring Lemma \cite{Kosaraju1994}] \label{lemma:path_coloring}
Let $G$ be a directed graph satisfying two conditions: (1) each vertex has an
indegree of at most 2, an outdegree of at most 2, and a total degree of at most
3, and (2) the graph contains no 2-cycles. Under these conditions, the edges of
$G$ can be 2-colored such that each color class forms a collection of directed
paths.
\end{lemma}

\paragraph{Algorithm Analysis}
We introduce four weight quantities describing how the edges of $T_{opt}$
interact with the 2-cycles of $\mathcal{C}$:
\begin{itemize}
    \item $b_1 W$ -- the total weight of the edges in $T_{opt}$ that are the
    heavier edge of a 2-cycle in $\mathcal{C}$
    \item $c_1 W$ -- the total weight of the corresponding (lighter) sibling
    edges of those same 2-cycles
    \item $c_2 W$ -- the total weight of the edges in $T_{opt}$ that are the
    lighter edge of a 2-cycle in $\mathcal{C}$
    \item $b_2 W$ -- the total weight of the corresponding (heavier) sibling
    edges of those same 2-cycles
\end{itemize}

The weight of the optimal tour edges not belonging to a 2-cycle in $\mathcal{C}$
is $w(T_{opt}) - b_1W - c_2W$. By adding back the modified weight of the 2-cycle
edges, the total weight of the optimal tour in $G'$ is bounded below by:
$$w(T_{opt}) - b_1W - c_2W + 2(b_1 - c_1)W + 2(b_2 - c_2)W$$

Given that each 2-cycle $C_i$ satisfies $c_i \le b_i \le 2c_i$, it follows that:
$$(b_1 - 2c_1 + 2b_2 - 3c_2) \ge (b_1 + b_2) - 2(c_1 + c_2) \ge b - 2c$$
This implies the optimal tour's weight in $G'$ is at least $w(T_{opt}) + (b -
2c)W$. As a result, the matching extracted in Step 2b has a weight of at least
$\frac{1}{2}(w(T_{opt}) + (b - 2c)W)$.

In Step 2c, including the $3^+$-cycles adds a weight of $(1 - p - b - c)W$.
Therefore, prior to Step 2d, the graph's total weight is at least:
\begin{equation} \label{eq:weight_before_contraction_real}
\frac{w(T_{opt}) + (b - 2c)W}{2} + (1 - p - b - c)W
\end{equation}

To evaluate the remaining weight after Step 2d, let $A$ denote the set of
2-cycles with corresponding edges in the matching. Let $\hat{b}W$ and $\hat{c}W$
be the sum of the weights of the heavier and lighter edges in $A$, respectively.
Simplifying equation (\ref{eq:weight_before_contraction_real}) and subtracting
the contracted weight $2(\hat{b} - \hat{c})W$, we are left with an edge
collection at the end of Step 2d weighing at least:
\begin{equation} \label{eq:weight_after_contraction_real}
\frac{w(T_{opt})}{2} + \left(1 - 2c - \frac{b}{2} - p - 2(\hat{b} - \hat{c})\right)W
\end{equation}

\paragraph{Handling Induced 2-Cycles}
The edges at the start of Step 2e consist of paths originating from the
$3^+$-cycles combined with matching edges. While this graph violates the
conditions of Lemma \ref{lemma:path_coloring}, the following lemma proves it can
still be path-colored with two colors.

\begin{lemma}[\cite{Kosaraju1994}] \label{lemma:induced_cycles}
The set of edges $E$ remaining after Step 2d can be successfully path-colored
using two colors.
\end{lemma}

\begin{proof}
By design, every vertex has at most one incoming cycle edge, one outgoing cycle
edge, and one matching edge, satisfying the degree constraints of Lemma
\ref{lemma:path_coloring}. Moreover, contracting the 2-cycles of $\mathcal{C}$
merges two vertices (each of degree at most 3, sharing 2 edges) into a single
vertex of degree at most 2. This satisfies all conditions of Lemma
\ref{lemma:path_coloring} except for the potential presence of 2-cycles. These
remaining 2-cycles, formed solely by one cycle edge and one matching edge, are
referred to as \textit{induced 2-cycles}.

Consider such an induced 2-cycle $\{(u, v), (v, u)\}$. Due to the degree-three
limit and the presence of a matching edge between $u$ and $v$, there is at most
one additional edge connected to $v$ (labeled $(v, x)$) and at most one
additional edge connected to $u$ (labeled $(u, y)$), neither being a matching
edge. To resolve the 2-cycle, the path $(x, v, u, y)$ is contracted into a
single edge $(x, y)$. The resulting graph retains the same properties but
contains one fewer 2-cycle. This process is repeated until all 2-cycles are
eliminated. Once the contracted graph is colored, coloring the original graph is
straightforward: if $(x, y)$ is colored red, then $(x, v)$, $(v, u)$, and $(u,
y)$ are colored red, while $(u, v)$ is assigned blue. Because $(u, v)$ shares no
vertices with another blue edge, this assignment is perfectly safe.
\end{proof}

\paragraph{An Inconsistency Between the Two Sources}
The proof above treats induced 2-cycles as the only way in which the edges
entering Step 2e can violate the conditions of Lemma
\ref{lemma:path_coloring}. A matching edge may however also run in the
same direction as a kept cycle edge, and the two then form a double
edge. Kosaraju, Park and Stein implicitly accept such an input: their lemma is
stated for a directed graph without further restriction, and their case
analysis does not list double edges among the situations to be resolved.
Bl\"aser \cite{Blaser2004}, whose proof of that lemma we reconstruct in
Chapter \ref{chap:path_coloring}, reads it the other way: he states it for a
loopless graph and notes explicitly that a graph containing double edges does
not fulfil its premises. The two sources therefore disagree about what the
path-coloring lemma may be given.

Step 2e of the implementation therefore contracts double edges alongside the
2-cycles, so that the path coloring algorithm receives the input Bl\"aser's proof assumes.

\paragraph{Final Bound}
Given that the edge weight remaining after Step 2d's contractions is bounded by
equation (\ref{eq:weight_after_contraction_real}), the preceding lemma allows us
to secure a collection of directed paths with a weight of at least:
\begin{equation} \label{eq:weight_directed_paths_real}
\frac{1}{2} \left( \frac{w(T_{opt})}{2} + \left(1 - 2c - \frac{b}{2} - p - 2(\hat{b} - \hat{c})\right)W \right)
\end{equation}
The final step—adding 2-cycle edges at the conclusion of Step 2e—increases the
value in (\ref{eq:weight_directed_paths_real}) by a minimum of $(\hat{b} + (c -
\hat{c}))W$. This is because, for each 2-cycle in $\mathcal{C}$, we include its
heavier edge if the corresponding undirected edge was in $M$, or an edge
weighing at least as much as the lighter edge if it was not. Since $\mathcal{C}$
is a cycle cover, the vertices of a 2-cycle belong to no other cycle, so if its
collapsed edge was chosen in $M$, both endpoints are matched only to each other,
the contracted vertex is then completely isolated in the graph colored in Step
2e, and either directed edge can be added back without conflict, so we pick the
heavier one. Otherwise, at least one endpoint may carry an external matching
edge, which fixes which direction is safe to add -- we can then only guarantee
the weight of the lighter edge. Finally, this gives a path collection with a
weight of at least:
\begin{align*}
& \frac{1}{2} \left( \frac{w(T_{opt})}{2} + \left(1 - 2c - \frac{b}{2} - p - 2(\hat{b} - \hat{c})\right)W \right) + (\hat{b} + (c - \hat{c}))W \\
& \ge \left( \frac{3}{4} - \frac{b}{4} - \frac{p}{2} \right) w(T_{opt})
\ge \left( \frac{7}{12} - \frac{1}{12}(b - 2c) \right) w(T_{opt}).
\end{align*}

This final inequality is derived by applying the bound $p \le \frac{1}{3}(1 - b
- c)$, which remains valid even after reclassifying 2-cycles with $b_i > 2c_i$
as $3^+$-cycles. This leads to the final lemma.

\begin{theorem} \label{lemma:version2_final}
Version 2 of the generic algorithm guarantees a $\left(\frac{7}{12} -
\frac{1}{12}(b - 2c)\right)$-approximation for the Max TSP.
\end{theorem}

Although this bound is formulated using the values of $b$ and $c$ obtained after
reclassifying specific 2-cycles, this reclassification can only decrease the
value of $b - 2c$. Therefore, the bound established in Theorem
\ref{lemma:version2_final} is strictly valid for the original parameters of $b$
and $c$ as well.

\paragraph{Computational Complexity}
Let $n$ denote the number of vertices of $G$. Building the auxiliary graph $G'$
in Step 2a takes $O(n^2)$ time. Extracting the maximum-weight matching $M$ in
Step 2b, using Edmonds' blossom algorithm, takes $O(n^3)$ time
\cite{edmonds1965paths, edmonds1965matching}. Steps 2c and 2d are linear in $n$. The path coloring of
Step 2e executes deterministically in $O(n^3)$ time, as established in Chapter
4. The matching computation and path coloring therefore dominate, and Version 2
runs in $O(n^3)$ time overall.
